

\documentclass[12pt]{article}

% Packages
\usepackage[margin=1in]{geometry}
\usepackage{setspace}
\usepackage{titlesec}
\usepackage{parskip}

% Formatting
\setstretch{1.15}
\titleformat{\section}{\large\bfseries}{}{0em}{}

% Title information
\title{Peer Review: Martin Huang}
\author{Alejandro De La Torre \\
ECON 580 -- Honors Research Tutorial}
\date{\today}

\begin{document}

\maketitle

\section{Research Question and Motivation}

Martin Huang’s draft asks how temperature shocks affect labor supply in the United States over time, with a particular focus on whether deviations from normal temperatures reduce hours worked and whether those effects are temporary or persistent across months. This is a relevant and timely question because climate change is making extreme weather more frequent, and labor supply is one of the most direct channels through which weather may affect economic output. The draft also does a good job centering the United States, which is useful because the project can draw on substantial geographic variation in climate, occupation, and local labor-market structure. At the same time, the proposal is still quite broad. The next draft would benefit from being much more precise about the time period under study, what exactly counts as a temperature shock, and what is meant by labor supply in practice.


\section{Strengths}

This project has several clear strengths. First, the topic is important and grounded in a real economic issue with obvious policy relevance. Second, the motivation is intuitive: when temperatures become unusually hot or cold, some forms of work become harder, more dangerous, or even temporarily impossible. Third, the author is asking a question that could matter a great deal if properly narrowed, especially for understanding how climate risk affects working people in exposed sectors. I also think the choice to focus on the United States is a strength because it avoids becoming too globally diffuse and gives the paper a more coherent institutional setting. Finally, the paper already points toward a useful distinction between short-run disruption and repeated cumulative exposure, which could become a strong contribution if the author clearly identifies which workers are being affected and through what mechanism.


\section{Suggestions for Improvement}

My main suggestion is that the paper needs to get much more specific about who is actually being studied. Right now, ``labor supply'' is too broad. The author should define whether the outcome is weekly hours worked, monthly hours worked, labor force participation, employment, missed work days, or something else. Relatedly, the paper should clarify whose labor supply is likely to respond to temperature shocks. It seems unlikely that all workers are equally exposed. A white-collar worker in an insulated office, or someone working remotely, is very different from a construction worker, farmworker, delivery worker, mechanic, warehouse worker, or other worker whose job involves physical presence in heat or cold. Because of that, I would strongly encourage the author to narrow the analysis to occupations or sectors where the mechanism is strongest, or at least to build heterogeneity into the empirical design.

A second major issue is the definition of a temperature shock. The current draft refers to monthly temperature data and deviations from normal temperature levels, but this still feels vague. A shock could mean an unusually hot day, a multi-day heat wave, a cold snap, or a monthly average deviation from a local historical baseline. These are not the same thing. I think the author should explain whether the relevant shock is defined relative to each state’s own climate, whether extreme heat and extreme cold are both included, and whether duration matters as much as intensity. This is especially important because a nominal temperature increase may mean very different things in Arizona, Wisconsin, and coastal California. A state that is consistently hot may not experience the same kind of labor-market disruption from a small increase as a place where that temperature is highly unusual.

Third, I think the identification strategy needs much more development. As written, the project risks staying at the level of broad association. If the goal is to say something causal, the paper needs to explain why variation in temperature shocks can be interpreted as generating meaningful differences in labor outcomes rather than simply moving alongside other labor-market changes. This is especially important because labor supply can be affected by many other factors at the same time: COVID, local recessions, infrastructure disruptions, school closures, migration flows, industry composition changes, and commuting constraints, among others. If one month shows a relationship between a weather shock and lower hours worked but a similar month does not, that may indicate that non-weather forces are also driving the result. I would therefore encourage the author to think hard about what controls or fixed effects would be necessary, and about whether a more local unit of analysis would help.

Fourth, the paper should think more seriously about geography and exposure. State-level monthly averages may be too coarse if actual exposure varies meaningfully within states. Temperatures can differ across counties, commuting zones, and even smaller areas, and workers may live in one place and work in another. This creates potential measurement error if the paper assigns a single temperature to everyone in a state. If more granular weather data are feasible, that would likely improve the design. If not, then the paper should at least acknowledge the limitation and justify why state-level analysis is still informative.

Fifth, I think the project should be more open to heterogeneous and even offsetting effects. The current motivation reads somewhat as if temperature shocks should lower labor supply on average, but that may not always be true. In some sectors, unusually hot weather might reduce hours; in others, it may increase labor demand. For example, repair work, cooling-related services, seasonal recreation, or climate adaptation work may expand during certain shocks. Likewise, broader climate-related events may reduce employment in one sector while increasing it in another. If the analysis only looks at an aggregate labor-supply measure, the effects could wash out even when the shock is having real consequences for particular workers. This is one reason why I think disaggregating by occupation, industry, gender, rural versus urban setting, or other relevant dimensions would make the project much stronger.

Relatedly, I think the author should consider the possibility of labor substitution or compositional change rather than only reduced labor supply among the same workers. In some places, a climate-related disruption may push one kind of labor out while drawing another kind in. For instance, if one set of jobs becomes less viable because of repeated climate exposure, there may simultaneously be new labor demand for repair, mitigation, rebuilding, or adaptation projects. In that case, an aggregate measure might show little overall change even though there is substantial turnover underneath. More generally, migration into and out of states or metro areas could also affect the measured labor-supply response if incoming workers are concentrated in occupations that are relatively insulated from weather. For that reason, it would help the paper to think about population flows, sectoral composition, and the distinction between a change in hours for incumbent workers versus a change in who is working.

Finally, I would recommend tightening the literature review. The papers listed are relevant, but the next draft should identify the few papers that are most central to this exact question, especially papers on labor supply, hours worked, weather exposure, and adaptation in the United States. That would help the author clarify what gap this project is actually filling and whether the contribution is primarily about the United States, about persistence over time, or about a particular type of worker.


\section{Minor Comments}

A few smaller comments may also help strengthen the draft. The paper should clarify whether it is really about temperature alone or about broader weather shocks more generally, since storms, humidity, snow, hurricanes, and other prolonged weather events may matter as much as temperature in some places. The proposal would also benefit from explicitly naming likely data sources for both labor outcomes and weather measures so the reader can better judge feasibility. More broadly, I think this is a strong and highly relevant idea. The key now is to ground it in a more clearly defined population, outcome, and empirical design so that the paper can say something meaningful about which workers are actually affected and under what conditions.

\end{document}